% ---
% RESUMOS
% ---

% resumo em português
\setlength{\absparsep}{18pt} % ajusta o espaçamento dos parágrafos do resumo
\begin{resumo}
O controle patrimonial em ambientes acadêmicos pode ser dificultado pela movimentação de equipamentos, pela dependência de conferência manual e pela ocorrência de divergências entre registros administrativos e localização física dos bens. Este trabalho apresenta o desenvolvimento de um protótipo funcional de sistema web para inventário patrimonial baseado em Identificação por Radiofrequência (RFID), aplicado ao contexto do Colegiado de Ciência da Computação da Universidade Estadual de Santa Cruz. A solução integra interface web, API REST e processamento de eventos RFID para apoiar o acompanhamento dos bens, a auditoria patrimonial, o registro histórico e a identificação de inconsistências. A arquitetura preserva o fluxo previsto para sensores ou leitores em rede, enquanto a validação física concentrou-se em um cenário controlado com leitor RFID de proximidade de baixo custo e uma tag. Os resultados indicam que o protótipo valida o núcleo funcional da proposta, pois recebe eventos de leitura, processa tags, atualiza informações do inventário e registra divergências. Ensaios com maior alcance, múltiplas tags e antenas RFID de longo alcance permanecem como continuidade experimental.

 \textbf{Palavras-chave}: Inventário patrimonial. RFID. Sistema web. Auditoria patrimonial. Rastreabilidade.
\end{resumo}

% resumo em inglês
\begin{resumo}[Abstract]
 \begin{otherlanguage*}{english}
The control of assets in academic environments can be hindered by equipment movement, reliance on manual checking, and inconsistencies between administrative records and the physical location of assets. This work presents the development of a functional web system prototype for asset inventory based on Radio-Frequency Identification (RFID), applied to the context of the Computer Science Department at the State University of Santa Cruz. The solution integrates a web interface, a REST API, and RFID event processing to support asset monitoring, asset auditing, operational history, and inconsistency identification. The architecture preserves the flow planned for networked sensors or readers, while physical validation focused on a controlled scenario using a low-cost short-range RFID reader and a single tag. The results indicate that the prototype validates the functional core of the proposal, since it receives reading events, processes tags, updates inventory information, and records divergences. Experiments with greater range, multiple tags, and long-range RFID antennas remain as future work.

   \vspace{\onelineskip}
 
   \noindent 
   \textbf{Keywords}: Asset inventory. RFID. Web system. Asset auditing. Traceability.
 \end{otherlanguage*}
\end{resumo}

